% GENERATED FILE. Edit canonical lessons, then run npm run generate:books.
% canonical-source-hash: fnv1a64:e326989a75947e0d
% canonical-lessons: ML-C480-yatravivaranam, ML-C480-marannuvaykkuka, ML-C480-codiccariyuka, ML-C480-lokkar, ML-C480-karutivaykkuka

\chapter{Keeping Things Safe}
\label{ch:ml-keeping-things-safe}
\hlchaptermodality{480}

\begin{chapteropening}
\textbf{By the end of this chapter:} \emph{I can say a travelogue, to leave behind, to ask around, a locker and to keep aside.}

Complete the last lesson of chapter 480: I can say a travelogue, to leave behind, to ask around, a locker and to keep aside.
\end{chapteropening}

\section[\texorpdfstring{\ml{യാത്രാവിവരണം} (yātrāvivaraṇaṁ)}{യാത്രാവിവരണം (yātrāvivaraṇaṁ)}]{\ml{യാത്രാവിവരണം} (yātrāvivaraṇaṁ) — a travelogue}

\label{lesson:ML-C480-yatravivaranam}



\begin{quote}
Before the new one: say the Malayalam for a stopover, then the Malayalam for to pass by.
\end{quote}

\subsection*{You'll want to know: \ml{യാത്രാവിവരണം}}
\textbf{\ml{യാത്രാവിവരണം}} — \emph{yātrāvivaraṇaṁ} — \textquotedblleft{}a travelogue\textquotedblright{}.

\begin{quote}
\ml{അവൻ} \ml{തന്റെ} \ml{യാത്രയെക്കുറിച്ച്} \ml{ഒരു} \ml{യാത്രാവിവരണം} \ml{എഴുതി}.
\end{quote}

\begin{quote}
\emph{avan tanṟe yātrayekkuṟiccŭ oru yātrāvivaraṇaṁ eḻuti.}
\end{quote}

\begin{quote}
He wrote a travelogue about his journey.
\end{quote}

\ml{വിവരണം}, an account, is the noun of \ml{വിവരിക്കുക}, to describe.

\begin{cousinweb}
\textbf{\ml{യാത്രാവിവരണം}} (\emph{yātrāvivaraṇaṁ}) is a Malayalam compound built from \ml{യാത്ര}, from Sanskrit \emph{yātrā} (CDIAL 10456), and \ml{വിവരണം}, an account, from Sanskrit \emph{vivaraṇa}, \textquotedblleft{}laying open\textquotedblright{} (Monier-Williams).
\end{cousinweb}

\begin{grammarlens}[title={What you've built}]
One more word for getting around.
\end{grammarlens}

\subsection*{Your turn}
\begin{itemize}
\raggedright
  \item \emph{Say it:} \emph{yātrāvivaraṇaṁ}
  \item \emph{Say it:} \emph{yātrāvivaraṇaṁ}, once more
  \item \emph{Say it:} \emph{avan tanṟe yātrayekkuṟiccŭ oru yātrāvivaraṇaṁ eḻuti.} — He wrote a travelogue about his journey.
  \item \emph{Recall:} say the Malayalam for a stopover, then the Malayalam for to pass by, then say \emph{yātrāvivaraṇaṁ} again
\end{itemize}

\begin{tcolorbox}[breakable,colback=teal!4,colframe=teal!35!black,title={Before you move on}]
What does \ml{യാത്രാവിവരണം} mean? (\textquotedblleft{}A travelogue\textquotedblright{}.) Say it once more.
\end{tcolorbox}
\section[\texorpdfstring{\ml{മറന്നുവയ്ക്കുക} (maṟannuvaykkuka)}{മറന്നുവയ്ക്കുക (maṟannuvaykkuka)}]{\ml{മറന്നുവയ്ക്കുക} (maṟannuvaykkuka) — to leave behind (by forgetting)}

\label{lesson:ML-C480-marannuvaykkuka}



\begin{quote}
Before the new one: say the Malayalam for to pass by, then the Malayalam for a travelogue.
\end{quote}

\subsection*{You'll want to know: \ml{മറന്നുവയ്ക്കുക}}
\textbf{\ml{മറന്നുവയ്ക്കുക}} — \emph{maṟannuvaykkuka} — \textquotedblleft{}to leave behind (by forgetting)\textquotedblright{}.

\begin{quote}
\ml{ഞാൻ} \ml{എന്റെ} \ml{കുട} \ml{ബസിൽ} \ml{മറന്നുവച്ചു}.
\end{quote}

\begin{quote}
\emph{ñān enṟe kuṭa basil maṟannuvaccu.}
\end{quote}

\begin{quote}
I left my umbrella on the bus.
\end{quote}

Past: \ml{മറന്നുവച്ചു}, \emph{maṟannuvaccu}.

Literally to forget and put: \ml{മറന്ന്}, having forgotten, plus \ml{വയ്ക്കുക}, to put down.

\begin{cousinweb}
\textbf{\ml{മറന്നുവയ്ക്കുക}} (\emph{maṟannuvaykkuka}) is a native Dravidian compound of \ml{മറക്കുക}, to forget (DEDR 4760), and \ml{വയ്ക്കുക}, to put (DEDR 5549).
\end{cousinweb}

\begin{grammarlens}[title={What you've built}]
One more word for getting around.
\end{grammarlens}

\subsection*{Your turn}
\begin{itemize}
\raggedright
  \item \emph{Say it:} \emph{maṟannuvaykkuka}
  \item \emph{Say it:} \emph{maṟannuvaykkuka}, once more
  \item \emph{Say it:} \emph{ñān enṟe kuṭa basil maṟannuvaccu.} — I left my umbrella on the bus.
  \item \emph{Recall:} say the Malayalam for to pass by, then the Malayalam for a travelogue, then say \emph{maṟannuvaykkuka} again
\end{itemize}

\begin{tcolorbox}[breakable,colback=teal!4,colframe=teal!35!black,title={Before you move on}]
What does \ml{മറന്നുവയ്ക്കുക} mean? (\textquotedblleft{}To leave behind (by forgetting)\textquotedblright{}.) Say it once more.
\end{tcolorbox}
\section[\texorpdfstring{\ml{ചോദിച്ചറിയുക} (cōdiccaṟiyuka)}{ചോദിച്ചറിയുക (cōdiccaṟiyuka)}]{\ml{ചോദിച്ചറിയുക} (cōdiccaṟiyuka) — to ask around; to find out by asking}

\label{lesson:ML-C480-codiccariyuka}



\begin{quote}
Before the new one: say the Malayalam for a travelogue, then the Malayalam for to leave behind.
\end{quote}

\subsection*{You'll want to know: \ml{ചോദിച്ചറിയുക}}
\textbf{\ml{ചോദിച്ചറിയുക}} — \emph{cōdiccaṟiyuka} — \textquotedblleft{}to ask around; to find out by asking\textquotedblright{}.

\begin{quote}
\ml{ഞാൻ} \ml{ഒരു} \ml{കടക്കാരനോട്} \ml{വഴി} \ml{ചോദിച്ചറിഞ്ഞു}.
\end{quote}

\begin{quote}
\emph{ñān oru kaṭakkāranōṭŭ vaḻi cōdiccaṟiññu.}
\end{quote}

\begin{quote}
I found out the way by asking a shopkeeper.
\end{quote}

Past: \ml{ചോദിച്ചറിഞ്ഞു}, \emph{cōdiccaṟiññu}.

\ml{ചോദിച്ച്}, having asked, plus \ml{അറിയുക}, to know. The person asked takes -\ml{ഓട്}.

\begin{cousinweb}
\textbf{\ml{ചോദിച്ചറിയുക}} (\emph{cōdiccaṟiyuka}) is a Malayalam compound built from \ml{ചോദിക്കുക}, to ask, and the native verb \ml{അറിയുക}, to know (DEDR 314).
\end{cousinweb}

\begin{grammarlens}[title={What you've built}]
One more word for getting around.
\end{grammarlens}

\subsection*{Your turn}
\begin{itemize}
\raggedright
  \item \emph{Say it:} \emph{cōdiccaṟiyuka}
  \item \emph{Say it:} \emph{cōdiccaṟiyuka}, once more
  \item \emph{Say it:} \emph{ñān oru kaṭakkāranōṭŭ vaḻi cōdiccaṟiññu.} — I found out the way by asking a shopkeeper.
  \item \emph{Recall:} say the Malayalam for a travelogue, then the Malayalam for to leave behind, then say \emph{cōdiccaṟiyuka} again
\end{itemize}

\begin{tcolorbox}[breakable,colback=teal!4,colframe=teal!35!black,title={Before you move on}]
What does \ml{ചോദിച്ചറിയുക} mean? (\textquotedblleft{}To ask around; to find out by asking\textquotedblright{}.) Say it once more.
\end{tcolorbox}
\section[\texorpdfstring{\ml{ലോക്കർ} (lōkkar)}{ലോക്കർ (lōkkar)}]{\ml{ലോക്കർ} (lōkkar) — a locker}

\label{lesson:ML-C480-lokkar}



\begin{quote}
Before the new one: say the Malayalam for to leave behind, then the Malayalam for to ask around.
\end{quote}

\subsection*{You'll want to know: \ml{ലോക്കർ}}
\textbf{\ml{ലോക്കർ}} — \emph{lōkkar} — \textquotedblleft{}a locker\textquotedblright{}.

\begin{quote}
\ml{ഹോട്ടലിലെ} \ml{ലോക്കറിൽ} \ml{ഞാൻ} \ml{പാസ്പോർട്ട്} \ml{വച്ചു}.
\end{quote}

\begin{quote}
\emph{hōṭṭalile lōkkaṟil ñān pāspōrṭṭŭ vaccu.}
\end{quote}

\begin{quote}
I put my passport in the hotel locker.
\end{quote}

Banks also rent out \ml{ലോക്കർ} space for gold and papers.

\begin{cousinweb}
\textbf{\ml{ലോക്കർ}} (\emph{lōkkar}) was borrowed from English \emph{locker}, a small cupboard that locks.
\end{cousinweb}

\begin{grammarlens}[title={What you've built}]
One more word for getting around.
\end{grammarlens}

\subsection*{Your turn}
\begin{itemize}
\raggedright
  \item \emph{Say it:} \emph{lōkkar}
  \item \emph{Say it:} \emph{lōkkar}, once more
  \item \emph{Say it:} \emph{hōṭṭalile lōkkaṟil ñān pāspōrṭṭŭ vaccu.} — I put my passport in the hotel locker.
  \item \emph{Recall:} say the Malayalam for to leave behind, then the Malayalam for to ask around, then say \emph{lōkkar} again
\end{itemize}

\begin{tcolorbox}[breakable,colback=teal!4,colframe=teal!35!black,title={Before you move on}]
What does \ml{ലോക്കർ} mean? (\textquotedblleft{}A locker\textquotedblright{}.) Say it once more.
\end{tcolorbox}
\section[\texorpdfstring{\ml{കരുതിവയ്ക്കുക} (karutivaykkuka)}{കരുതിവയ്ക്കുക (karutivaykkuka)}]{\ml{കരുതിവയ്ക്കുക} (karutivaykkuka) — to keep aside; to keep ready}

\label{lesson:ML-C480-karutivaykkuka}



\begin{quote}
Before the new one: say the Malayalam for to ask around, then the Malayalam for a locker.
\end{quote}

\subsection*{You'll want to know: \ml{കരുതിവയ്ക്കുക}}
\textbf{\ml{കരുതിവയ്ക്കുക}} — \emph{karutivaykkuka} — \textquotedblleft{}to keep aside; to keep ready\textquotedblright{}.

\begin{quote}
\ml{യാത്രയ്ക്കായി} \ml{അമ്മ} \ml{കുറച്ച്} \ml{പണം} \ml{കരുതിവച്ചു}.
\end{quote}

\begin{quote}
\emph{yātraykkāyi amma kuṟaccŭ paṇaṁ karutivaccu.}
\end{quote}

\begin{quote}
Mother kept some money aside for the journey.
\end{quote}

Past: \ml{കരുതിവച്ചു}, \emph{karutivaccu}.

\ml{കരുതി}, keeping in mind, plus \ml{വയ്ക്കുക}, to put: to put something by for later.

\begin{cousinweb}
\textbf{\ml{കരുതിവയ്ക്കുക}} (\emph{karutivaykkuka}) is a Malayalam compound built from \ml{കരുതുക}, to keep in mind, and the native verb \ml{വയ്ക്കുക}, to put (DEDR 5549).
\end{cousinweb}

\begin{grammarlens}[title={What you've built}]
One more word for getting around.
\end{grammarlens}

\subsection*{Your turn}
\begin{itemize}
\raggedright
  \item \emph{Say it:} \emph{karutivaykkuka}
  \item \emph{Say it:} \emph{karutivaykkuka}, once more
  \item \emph{Say it:} \emph{yātraykkāyi amma kuṟaccŭ paṇaṁ karutivaccu.} — Mother kept some money aside for the journey.
  \item \emph{Recall:} say the Malayalam for to ask around, then the Malayalam for a locker, then say \emph{karutivaykkuka} again
\end{itemize}

\begin{tcolorbox}[breakable,colback=teal!4,colframe=teal!35!black,title={Before you move on}]
What does \ml{കരുതിവയ്ക്കുക} mean? (\textquotedblleft{}To keep aside; to keep ready\textquotedblright{}.) Say it once more.
\end{tcolorbox}
